\documentclass{beamer}

\mode<presentation>





\usetheme{Frankfurt}%
\usecolortheme{seagull}
\logo{\includegraphics[height=.25in]{clarksonGreen}}

\definecolor{garnet}{RGB}{136,0,0}
%\definecolor{clarksonGreen}{RGB}{0,71,28}
\definecolor{clarksonGreen}{RGB}{0,52,21}
\setbeamercolor{palette primary}{fg=clarksonGreen,bg=white}
\setbeamercolor{palette secondary}{fg=clarksonGreen,bg=white}
\setbeamercolor{palette tertiary}{fg=clarksonGreen,bg=white}
\setbeamercolor{palette quaternary}{bg=clarksonGreen,fg=white}
\setbeamercolor{block title}{fg=black,bg=black!15}
\setbeamercolor{block body}{fg=black,bg=black!10}
\setbeamercolor{titlelike}{bg=clarksonGreen,fg=white} % parent=palette quaternary}

\newcommand{\half}{\mbox{$\frac{1}{2}$}}
\newcommand{\deltat}{\mbox{$\triangle t$}}
\newcommand{\deltax}{\mbox{$\triangle x$}}
\newcommand{\deltay}{\mbox{$\triangle y$}}

\newcommand{\deriv}[2]{\frac{d}{d#2}#1}
\newcommand{\derivTwo}[2]{\frac{d^2}{d#2^2}#1}

\newcommand{\lp}{\left(}
\newcommand{\rp}{\right)}



\newcommand{\arrayTwo}[4]{
  \left[
  \begin{array}{rr}
    #1 & #2 \\
    #3 & #4
  \end{array}
  \right]
}


\newcommand{\vecTwo}[2]{
  \left[
  \begin{array}{r}
    #1 \\  #2
  \end{array}
  \right]
}


\newcommand{\vecThree}[3]{
  \left[
  \begin{array}{r}
    #1 \\  #2 \\ #3
  \end{array}
  \right]
}


\newcommand{\arrayThree}[9]{
  \left[
    \begin{array}{rrr}
      #1 & #2 & #3 \\
      #4 & #5 & #6 \\
      #7 & #8 & #9
    \end{array}
  \right]
}


\newcommand{\startRowFour}{
  \left[
    \begin{array}{rrr|r}
}

\newcommand{\oneRowFour}[4] {
      #1 & #2 & #3 & #4 \\
}

\newcommand{\stopRowOps}{
    \end{array}
  \right]
}




\begin{document}

\title{Ordinary Differential Equations}
\subtitle{Math 232 - Week 11, Day 3}

\author{Kelly Black}
\institute{Clarkson University}
\date{11 November 2011}

\begin{frame}
  \titlepage
\end{frame}

\begin{frame}
  \frametitle{Outline}
  \tableofcontents[pausesection,hideallsubsections]
\end{frame}


\section{Complex Valued Eigen Values}


\begin{frame}
  \frametitle{Complex Valued Eigen Values}

  First, recall:
  \begin{eqnarray*}
    e^{a+bi} & = & e^a e^{bi}, \\
    & = & e^a \lp \cos(b) + i \sin(b) \rp.
  \end{eqnarray*}

  Also, the complex conjugate:
  \begin{eqnarray*}
    \overline{a+ib} & = & a-i b.
  \end{eqnarray*}

  Note: If A is a real valued matrix with complex valued eigen values,
  then the eigen values and eigen vectors come in complex conjugates
  pairs.

\end{frame}


\begin{frame}
  \frametitle{Example}

  \begin{eqnarray*}
    \deriv{~}{t} \vec{x} & = & \arrayTwo{5}{-10}{1}{-1} \vec{x}.
  \end{eqnarray*}

  \uncover<2->
  {
    Find the eigen values:
    \begin{eqnarray*}
      \det\lp\arrayTwo{5-\lambda}{-10}{1}{-1-\lambda}\rp & = & \lambda^2-4 \\
      \lambda_{1,2} & = & 2 \pm i.
    \end{eqnarray*}
  }

\end{frame}


\begin{frame}
  \frametitle{Find the Eigen Vectors}

  $\lambda_1 = 2+i$:
  \begin{eqnarray*}
    \arrayTwo{3-i}{-10}{1}{-3-i} \vecTwo{x}{y} & = & \vecTwo{0}{0} \\
    y & = & \lp \frac{3}{10} - \frac{i}{10} \rp x, \\
    \Rightarrow \vec{v}_1 & = & x \lp \vecTwo{1}{\frac{3}{10}} + \vecTwo{0}{\frac{-i}{10}} \rp.
  \end{eqnarray*}

  \uncover<2->
  {
    $\lambda_2 = 2-i$:
    \begin{eqnarray*}
      \arrayTwo{3+i}{-10}{1}{-3+i} \vecTwo{x}{y} & = & \vecTwo{0}{0} \\
      y & = & \lp \frac{3}{10} + \frac{i}{10} \rp x, \\
      \Rightarrow \vec{v}_2 & = & x \lp \vecTwo{1}{\frac{3}{10}} + \vecTwo{0}{\frac{i}{10}} \rp.
    \end{eqnarray*}
  }

\end{frame}


\begin{frame}
  \frametitle{Determine the Solution}

  \begin{eqnarray*}
    \uncover<1->
    {
      \vec{x} & = & A e^{(2+i)t} \lp \vecTwo{1}{\frac{3}{10}} + \vecTwo{0}{\frac{-i}{10}} \rp
      + B e^{(2-i)t} \lp \vecTwo{1}{\frac{3}{10}} - \vecTwo{0}{\frac{-i}{10}} \rp \\
    }
    \uncover<2->
    {
      \vec{x} & = & A e^{(2+i)t} \lp \underbrace{\vecTwo{1}{\frac{3}{10}}}_{\vec{p}} + 
      \underbrace{\vecTwo{0}{\frac{-i}{10}}}_{i\vec{q}} \rp
      + B e^{(2-i)t} \lp \underbrace{\vecTwo{1}{\frac{3}{10}}}_{\vec{p}} -
      \underbrace{\vecTwo{0}{\frac{-i}{10}}}_{i\vec{q}} \rp \\
    }
    \uncover<3->
    {
      & = & e^{2t} \lp
      A\lp \cos(t) + i \sin(t) \rp \lp \vec{p} + i \vec{q} \rp 
      + B \lp \cos(t) - i \sin(t) \rp \lp \vec{p} - i \vec{q} \rp 
      \rp
    }
  \end{eqnarray*}

\end{frame}


\begin{frame}
  \frametitle{More Algebra...}

  \begin{eqnarray*}
    \vec{x} & = & e^{2t} \lp
      A\lp \cos(t) + i \sin(t) \rp \lp \vec{p} + i \vec{q} \rp 
      + B \lp \cos(t) - i \sin(t) \rp \lp \vec{p} - i \vec{q} \rp 
      \rp \\
      & = & C_1 e^{2t} \lp \cos(t) \vec{p} - \sin(t) \vec{q} \rp
      + C_2 e^{2t} \lp \cos(t) \vec{q} + \sin(t) \vec{p} \rp.
  \end{eqnarray*}

  (The solutions spiral out.)

\end{frame}

\section{Characterization of Solutions}

\begin{frame}
  \frametitle{Characterization of Solutions}

  \begin{eqnarray*}
    \deriv{~}{t} \vec{x} & = & A \vec{x}.
  \end{eqnarray*}

  If $\lambda$'s are complex, $\lambda = a \pm ib$:
  \begin{enumerate}
  \item If $\alpha>0$ then it is a repelling spiral.
  \item If $\alpha<0$ then it is an attracting spiral.
  \item If $\alpha=0$ then the solutions are neutrally stable.
  \end{enumerate}

\end{frame}


\section{Nullclines}

\begin{frame}
  \frametitle{Nullclines}

  Given a solution to a differential equation,
  \begin{eqnarray*}
    \vec{x} & = & \vecTwo{x(t)}{y(t)},
  \end{eqnarray*}
  if
  \begin{itemize}
  \item $x'(t)=0$ then the trajectory is ``vertical'' at time $t$,
  \item $y'(t)=0$ then the trajectory is ``horizontal'' at time $t$.
  \end{itemize}

\end{frame}


\begin{frame}
  \frametitle{Definition of the Nullclines}

  The set of points, $\vecTwo{x}{y}$ where $x'(t)=0$ is called the v-nullcline. (Vertical slope.)

  The set of points, $\vecTwo{x}{y}$ where $y'(t)=0$ is called the h-nullcline. (Horizontal slope.)

\end{frame}


\begin{frame}
  \frametitle{Example}
  
  In the previous example we had
  \begin{eqnarray*}
    \deriv{~}{t} \vec{x} & = & \arrayTwo{5}{-10}{1}{-1} \vec{x}.
  \end{eqnarray*}

  \uncover<2->
  {
    Find the v-nullcline:
    \begin{eqnarray*}
      \deriv{x}{t} & = & 5x - 10y, \\
      0 & = & 5x - 10y, \\
      y & = & \frac{1}{2} x.
    \end{eqnarray*}
  }

  \uncover<3->
  {
    Find the h-nullcline:
    \begin{eqnarray*}
      \deriv{y}{t} & = & x - y, \\
      0 & = & x - y, \\
      y & = &  x.
    \end{eqnarray*}
  }


\end{frame}



\end{document}

% LocalWords:  Clarkson pausesection hideallsubsections
